\documentclass[a4paper]{article}
% Shared preamble for the English probability and statistics documents.
% Based on the original course preamble. Compile with pdfLaTeX.
\usepackage[top=1.5cm,bottom=2cm,left=2cm,right=2cm]{geometry}
\usepackage[T1]{fontenc}
\usepackage{lmodern}
\usepackage[utf8]{inputenc}
\usepackage[english]{babel}

\usepackage{amsmath}
\usepackage{amssymb}
\usepackage{amsthm}
\usepackage{mathrsfs}
\usepackage{mathtools}
\usepackage{etoolbox}
\usepackage{tikz}
\usetikzlibrary{shapes.multipart,arrows.meta}
\usepackage{xcolor}
\usepackage{enumerate}
\usepackage[shortlabels]{enumitem}
\usepackage{hyperref}
\usepackage{version}
\usepackage{bbold}
\usepackage{fontawesome5}
\usepackage{fancyhdr}
\lfoot{}
\usepackage{natbib}
\usepackage{graphicx}
\usepackage{pdfpages}
\usepackage{authblk}
    
\usepackage{blindtext}
\usepackage{multicol}
\usepackage{collcell}
\usepackage{wrapfig}
\usepackage{blkarray}
\usepackage{cancel}
\usepackage{appendix}
\usepackage{float}
\makeatletter
\input{tex/assets/slashbox.sty}
\makeatother
\hypersetup{unicode,colorlinks,citecolor=red,filecolor=black,
  linkcolor=red,urlcolor=cyan}
\setlength{\columnsep}{1cm}
\pagenumbering{gobble}

\def\numtext#1{%
  \ifcase#1\or one\or two\or three\or four\or five\or six\or seven\or eight\or nine\or ten\or
  eleven\or twelve\or thirteen\or fourteen\or fifteen\or sixteen\or seventeen\or eighteen\or nineteen\or twenty\or Lots\fi}

\newcommand{\Byte}[1]{%
  \def\nodecontents{}%
  \foreach[count=\l] \k in {#1} {\xappto{\nodecontents}{\noexpand\nodepart{\numtext{\l}}\k}}%
  \node[rectangle split,rectangle split parts=8,rectangle split horizontal,draw] at (2,2){\nodecontents};}

\newcommand{\D}{{\mathrm{d}}}
\newcommand{\cE}{{\mathcal{E}}}
\newcommand{\cG}{{\mathcal{G}}}
\newcommand{\cX}{{\mathcal{X}}}
\newcommand{\cY}{{\mathcal{Y}}}
\newcommand{\cB}{{\mathcal{B}}}
\newcommand{\cF}{{\mathcal{F}}}
\newcommand{\cO}{{\mathcal{O}}}
\newcommand{\cS}{{\mathcal{S}}}
\newcommand{\cT}{{\mathcal{T}}}
\newcommand{\given}{\, \big| \,}
\newcommand{\cL}{{\mathscr{L}}}
\newcommand{\cD}{{\mathcal{D}}}
\newcommand{\cA}{{\mathscr{A}}}
\newcommand{\cP}{{\mathscr{P}}}
\newcommand{\cR}{{\mathscr{R}}}
\newcommand{\A}{{\mathcal A}}
\newcommand{\one}{\mathbb{1}}
\newcommand{\mathbbm}[1]{\mathbb{#1}}

\providecommand{\C}{}
\renewcommand{\C}{\mathbb C}
\newcommand{\R}{\mathbb R}
\newcommand{\Q}{\mathbb Q}
\newcommand{\N}{\mathbb N}
\newcommand{\T}{\mathbb T}
\newcommand{\F}{\mathbb F}
\renewcommand{\P}{\mathbb P}
\newcommand{\Z}{\mathbb Z}
\newcommand{\E}{\mathbb {E}}
\newcommand{\Pbb}[1]{\P\left( #1\right)}
\newcommand{\Ebb}[1]{\E\left[ #1\right]}

\newcommand{\lb}{\left (}
\newcommand{\rb}{\right )}
\newcommand{\lbb}{\left [}
\newcommand{\rbb}{\right ]}
\newcommand{\labs}{\left |}
\newcommand{\rabs}{\right |}
\newcommand{\lbrb}[1]{\lb #1 \rb}
\newcommand{\lbbrbb}[1]{\lbb#1\rbb}
\newcommand{\labsrabs}[1]{\labs#1\rabs}
\newcommand{\lbbrb}[1]{\lbb#1\rb}
\newcommand{\lbrbb}[1]{\lb#1\rbb}
\newcommand{\langlerangle}[1]{\langle#1\rangle}
\newcommand{\lbcurly}{\{}
\newcommand{\rbcurly}{\}}
\newcommand{\curly}[1]{\lbcurly#1\rbcurly}
\newcommand{\abs}{\labsrabs}
\newcommand{\ind}[1]{1_{#1}}
\newcommand{\indep}{\perp \!\!\! \perp}
\renewcommand{\d}{\, \mathrm{d}}

\newtheorem{theorem}{Theorem}
\newtheorem{problemCnt}{Lemma}
\theoremstyle{definition}
\newtheorem{problem}[problemCnt]{Problem}
\newtheorem{problembook}[problemCnt]{\faBook \,  Problem}
\newtheorem{problempencil}[problemCnt]{\faPencil* \,  Problem}
\newtheorem{problemstar}[problemCnt]{\faStar\,  Problem}
\newtheorem{problemstarpencil}[problemCnt]{\faStar \, \faPencil* \,   Problem}
\newtheorem{probleminteresting}[problemCnt]{\faSmile[regular] \,   Problem}
\newtheorem{solution}{Solution}

\newcommand{\fwcell}[1]{\makebox[\arraystretch\normalbaselineskip]{$#1$}}
\newenvironment{enumeratei}{\begin{enumerate}[(i)]}{\end{enumerate}}
\newenvironment{enumeratea}{\begin{enumerate}[(a)]}{\end{enumerate}}
\newenvironment{enumeraten}{\begin{enumerate}[\upshape 1.]}{\end{enumerate}}

% Compatibility helpers needed by a few files in the combined archive.
\AtBeginDocument{\noindent{\footnotesize English translation prepared with ChatGPT (OpenAI, GPT-6), 26 October 2026, for translation of the Bulgarian text into English and LaTeX formatting.
The Bulgarian originals remain the reference; translation may contain errors.}\par\smallskip}
\newcommand{\MissingGraphic}[1]{%
  \begin{center}\fbox{\parbox{0.82\linewidth}{\centering Missing source image: \texttt{#1}}}\end{center}}



\setlength{\emergencystretch}{3em}
\pagenumbering{arabic}
\begin{document}
\begin{center}\Large Probability and Statistics, Computer Science\par\large Test 2, 12 June 2022: solutions\end{center}
\noindent\textit{English transcription of the supplied handwritten solution notes.}
\section*{Problem 1}
\subsection*{1. Normalisation and moments}
\begin{align*}
1&=\int_0^1\int_0^{1-x}(cx+1)\,dy\,dx
=\int_0^1(cx+1)(1-x)\,dx\\
&=\frac c2+1-\frac c3-\frac12=\frac c6+\frac12,
\end{align*}
so $c=3$. Also
\begin{align*}
\E[XY]&=\int_0^1\int_0^{1-x}xy(cx+1)\,dy\,dx
=\int_0^1x(cx+1)\frac{(1-x)^2}{2}\,dx\\
&=\frac c2\left(\frac13-\frac24+\frac15\right)+\frac12\left(\frac12-\frac23+\frac14\right)
=\frac c{60}+\frac1{24}=\frac{11}{120}.
\end{align*}
For $0\le x\le1$, $f_X(x)=(cx+1)(1-x)$, hence
\[\E X=\int_0^1(cx^2+x)(1-x)\,dx
=\frac c3+\frac12-\frac c4-\frac13=\frac c{12}+\frac16=\frac5{12}.\]
Similarly,
\[\E X^2=\int_0^1(cx^3+x^2-cx^4-x^3)\,dx
=c\left(\frac14-\frac15\right)+\frac13-\frac14=\frac c{20}+\frac1{12}=\frac7{30},\]
\[DX=\frac7{30}-\left(\frac5{12}\right)^2=\frac{43}{720}\approx0.0597.\]
For $0\le y\le1$,
\[f_Y(y)=\int_0^{1-y}(cx+1)\,dx=\frac c2(1-y)^2+1-y.\]
Therefore
\begin{align*}
\E Y&=\frac c2\int_0^1y(1-y)^2\,dy+\int_0^1y(1-y)\,dy\\
&=\frac c2\left(\frac12-\frac23+\frac14\right)+\frac12-\frac13
=\frac c{24}+\frac16=\frac7{24},\\
\E Y^2&=\frac c2\left(\frac13-\frac24+\frac15\right)+\frac13-\frac14
=\frac c{60}+\frac1{12}=\frac2{15},\\
DY&=\frac2{15}-\left(\frac7{24}\right)^2=\frac{139}{2880}\approx0.0483.
\end{align*}
Consequently,
\[Cor(X,Y)=\frac{\frac{11}{120}-\frac5{12}\frac7{24}}{\sqrt{\frac{43}{720}\frac{139}{2880}}}\approx-0.562,\qquad Cov(X,Y)=-\frac{43}{1440}\approx-0.030.\]
\subsection*{2. Conditional expectation}
For $0\le x\le1/2$,
\[f_{X\mid Y}(x\mid1/2)=\frac{f_{X,Y}(x,1/2)}{f_Y(1/2)}=\frac{cx+1}{7/8}.\]
Thus
\[\E[X\mid Y=1/2]=\int_0^{1/2}x\frac{cx+1}{7/8}\,dx
=\frac87\left(\frac c{24}+\frac18\right)=\frac27\approx0.2857.\]
\section*{Problem 2}
\begin{enumerate}
\item From $\P(N(\mu,10^2)<250)=15\%$, the standard normal table gives a quantile approximately $-1.04$ and the boxed result $\mu\approx260.4$.
\par\smallskip\textit{Transcription note: the handwritten standardisation line uses a denominator $10^2$ and then writes $250+100\cdot1.04=260.4$. The displayed arithmetic is inconsistent; the boxed mean is reproduced unchanged.}
\item
\[\P(N(\mu,10^2)>280)=\P\left(N(0,1)>\frac{280-\mu}{10}\right)
\approx1-\Phi(1.96)=\Phi(-1.96)\approx0.025=2.5\%.\]
\item
\begin{align*}
\P(230\le N(250,\sigma^2)\le270)
&=1-2\P(N(250,\sigma^2)<230)\\
&=1-2\Phi(-20/\sigma).
\end{align*}
The notes use $\Phi(-20/\sigma)=0.015$, hence $-20/\sigma\approx-2.17$ and $\sigma\approx9.2166$.
\end{enumerate}
\section*{Problem 3}
\begin{enumerate}
\item Let $X_i$ be student $i$'s score on problem 1. Then $X_i\sim10\cdot Ber(8/10)$, so
\[\E X_i=8,\qquad DX_i=100\frac8{10}\frac2{10}=16.\]
With $S_n=X_1+\cdots+X_n$, the CLT gives $(S_n-8n)/(4\sqrt n)\xrightarrow{d}N(0,1)$. Thus
\[\P(S_{100}/100<8)=\P\left(\frac{S_{100}-800}{4\sqrt{100}}<0\right)\approx\P(N(0,1)<0)=\frac12.\]
The notes also suggest reasoning by symmetry and write $S_{100}\sim10\cdot Bin(100,8/10)$.
\item Let $Y_i$ be the total points earned by student $i$, the sum of independent scores
\[X_i^{(1)}\sim10Ber(8/10),\quad X_i^{(2)}\sim10Ber(7/10),\quad X_i^{(3)}\sim10Ber(6/10),\quad X_i^{(4)}\sim10Ber(4/10).\]
Then $\E Y_i=10(8/10+7/10+6/10+4/10)=25$ and
\[DY_i=100\frac{8\cdot2+7\cdot3+6\cdot4+4\cdot6}{100}=85.\]
The grade is $Z_i=2+Y_i/10$, so $\E Z_i=4.5$ and $DZ_i=DY_i/100=17/20$.
\item By the CLT,
\[\frac{\frac1n\sum_{i=1}^nZ_i-4.5}{\sqrt{17/20}/\sqrt n}\xrightarrow{d}N(0,1).\]
Therefore
\begin{align*}
\P\left(\frac{Z_1+\cdots+Z_{100}}{100}-4.5>\frac{4.5}{100}\right)
&=\P\left(\frac{\frac1{100}\sum Z_i-4.5}{\sqrt{17/2000}}>\frac{4.5/100}{\sqrt{17/2000}}\right)\\
&\approx\P(N(0,1)>0.49)\approx31.21\%.
\end{align*}
\end{enumerate}
\section*{Problem 4}
\begin{enumerate}
\item Let $X\sim Exp(1/100)$. The desired probability is
\begin{align*}
\P(X<200)+\P(X\in[300,400])
&=1-\P(X\ge200)+\P(X\ge300)-\P(X\ge400)\\
&=1-e^{-2}+e^{-3}-e^{-4}\approx89.61\%,
\end{align*}
using $\P(X>t)=e^{-\lambda t}$.
\item Ins 2 pays 200 leva with probability $\P(X>300)=e^{-3}$. Its expected profit is
\[e^{-3}(x-200)+(1-e^{-3})x=x-200/e^3.\]
Set this equal to 10, giving $x\approx19.96$ leva.
\item Ins 1 pays $x=10+200/e^3$ to Ins 2 and retains a claim cost $X$ if $X<300$, or $X-200$ if $X\ge300$. Its expected retained claim is
\[\E X-200\P(X\ge300)=100-200e^{-3}\approx90.04\text{ leva}.\]
If the policy price is $y$, require
\[y-\E[\text{total cost}]=10,\qquad y-(x+100-200/e^3)=10,\]
so $y=120$. The margin asks: ``Can we say the answer is 120 regardless of the claim-size distribution?!'' Without reinsurance, $y-\E[\text{claim}]=10$, hence $y=110$.
\end{enumerate}
\end{document}
